\section{Review Scope and Method}
\label{sec:scope}

I conducted a purposive narrative review and finalized the source set in
July~2026.  The review focuses on how contract-based and runtime-assurance
approaches represent, compose, monitor, enforce, or repair assumptions that
cross robot-software and inter-robot interfaces.  Sources were compared using
four questions:

\begin{enumerate}
  \item What is the primary object of reasoning: mission, component behavior,
    dynamical interaction, distributed state, or runtime authority?
  \item Which mechanism is supplied: specification, refinement, composition,
    synthesis, monitoring, testing, intervention, or repair?
  \item What capability, assurance, or deployment evidence is reported, and
    does it include a physical multi-robot system?
  \item Which semantic, physical, temporal, uncertainty, evidence,
    compatibility, lifecycle, and authority obligations remain implicit across
    interfaces?
\end{enumerate}

I used title, abstract, and full-text keyword searches over arXiv and
web-indexed IEEE, ACM, Springer, Elsevier, and institutional repositories.  Core
queries combined \emph{assume--guarantee} or \emph{contract-based} with
\emph{robot}, \emph{ROS}, \emph{multi-robot}, \emph{multi-agent},
\emph{mission specification}, \emph{coordinated autonomy}, \emph{mission
reliability}, \emph{probabilistic model checking}, \emph{runtime verification},
and \emph{runtime assurance}.  Backward and forward snowballing from
foundational contract work, coordinated-mission assurance, robotic mission
frameworks, ROS contract and monitoring systems, and recent cross-layer reviews
added sources whose titles used different terminology.

I included primary work that defines a contract, compositional-assurance,
runtime-monitoring, or runtime-assurance mechanism applicable to robotics,
cyber-physical systems, or interconnected agents; reports a relevant robotics
implementation; or defines a directly relevant communication standard.  I also
included surveys that organize multi-robot modelling, communication, learning,
safety, or embodied assurance.  I excluded general design-by-contract work
without an architectural or cyber-physical connection and planning papers that
did not expose an interface-assurance question.  Representative planning and
middleware sources were retained only where they clarify what an interface must
carry.

The final cited corpus contains 40 sources spanning primary research, surveys,
foundational work, and standards.  For each source, I recorded the protected
object, mechanism, reported evidence, deployment setting, and obligations left
to adjacent components.  I then grouped sources by their main focus.
The Appendix codes these classes conservatively: ``E'' requires a property to be
explicit and central in the representative mechanisms; ``P'' covers mixed,
conditional, or model-internal treatment; and ``--'' marks a property outside
the stated mechanism.  When sources within one class differ, the class receives
``P'' unless the table separates the mechanisms.

I distinguish \emph{capability evidence}, showing that a system performs a
task; \emph{assurance evidence}, supporting a stated property under declared
assumptions; and \emph{deployment evidence}, measuring behavior on a stated
platform, configuration, workload, or fault condition.  The search evolved
through snowballing, so PRISMA-style candidate and exclusion counts
were not retained.  I make no prevalence or effect-size claim.  The evidence map
in Appendix~\ref{app:evidence-map} records qualitative coverage and keeps each
claim within the cited mechanism's scope.
